% PAGES: 282-282
% Straightforward continuation of section 9.4's Lagrangian derivation; no
% environments open or close on this page. Ends mid-sentence ("Then, the
% Hessian") -- plain prose continuing into sec09_c2_5.tex ("of $F_{0,1}(w)$
% is"), not a LaTeX construct, so nothing needs to be left open explicitly.

Note that
\begin{equation}
\left(\Sigma_R^{-1}\right)^t \;=\; \Sigma_R^{-1},
\end{equation}
since $\Sigma_R$ is a symmetric matrix; see Lemma 1.9. Then, from (9.49) and (9.50),
it follows that
\begin{equation}
w_0^t \;=\; -\frac{1}{2\lambda_0}\left(\Sigma_R^{-1}\bar\mu\right)^t \;=\;
-\frac{1}{2\lambda_0}\bar\mu^t\left(\Sigma_R^{-1}\right)^t \;=\;
-\frac{1}{2\lambda_0}\bar\mu^t\Sigma_R^{-1}.
\end{equation}

By substituting the formulas (9.49) for $w_0$ and (9.51) for $w_0^t$ in (9.48), we
obtain that
\begin{align}
\sigma_P^2 &= \frac{1}{4\lambda_0^2}\left(\bar\mu^t\Sigma_R^{-1}\right)\Sigma_R
\left(\Sigma_R^{-1}\bar\mu\right) \;=\;
\frac{1}{4\lambda_0^2}\bar\mu^t\Sigma_R^{-1}\left(\Sigma_R\Sigma_R^{-1}\right)\bar\mu
\notag\\
&= \frac{1}{4\lambda_0^2}\bar\mu^t\Sigma_R^{-1}\bar\mu.
\end{align}

Note that $\bar\mu^t\Sigma_R^{-1}\bar\mu > 0$ since $\Sigma_R$ is a symmetric positive
definite matrix, and therefore $\Sigma_R^{-1}$ is also symmetric positive definite;
cf. Lemma 5.4.

From (9.52), we find that
\[
\lambda_0^2 \;=\; \frac{1}{4\sigma_P^2}\bar\mu^t\Sigma_R^{-1}\bar\mu,
\]
which has two solutions, i.e.,
\begin{equation}
\lambda_{0,1} \;=\; -\frac{1}{2\sigma_P}\sqrt{\bar\mu^t\Sigma_R^{-1}\bar\mu};
\end{equation}
\begin{equation}
\lambda_{0,2} \;=\; \frac{1}{2\sigma_P}\sqrt{\bar\mu^t\Sigma_R^{-1}\bar\mu}.
\end{equation}

From (9.49), we find that the corresponding values of $w_0$ are
\begin{equation}
w_{0,1} \;=\; \frac{\sigma_P}{\sqrt{\bar\mu^t\Sigma_R^{-1}\bar\mu}}\Sigma_R^{-1}\bar\mu;
\end{equation}
\begin{equation}
w_{0,2} \;=\; -\frac{\sigma_P}{\sqrt{\bar\mu^t\Sigma_R^{-1}\bar\mu}}\Sigma_R^{-1}\bar\mu.
\end{equation}

Thus, the Lagrangian associated to this problem has two critical points,
$(w_{0,1},\lambda_{0,1})$ and $(w_{0,2},\lambda_{0,2})$.

To classify the critical point $(w_{0,1},\lambda_{0,1})$, let $F_{0,1}:\R^n\to\R$
given by $F_{0,1}(w)=F(w,\lambda_{0,1})$, where $\lambda_{0,1}$ is given by (9.53),
i.e.,
\begin{align*}
F_{0,1}(w) &= r_f+\bar\mu^tw+\lambda_{0,1}(w^t\Sigma_Rw-\sigma_P^2) \\
&= r_f+\bar\mu^tw-\frac{1}{2\sigma_P}\sqrt{\bar\mu^t\Sigma_R^{-1}\bar\mu}\,
(w^t\Sigma_Rw-\sigma_P^2).
\end{align*}

Note that $D^2(w^t\Sigma_Rw)=2\Sigma_R$, since $\Sigma_R$ is a symmetric matrix, and
$D^2(\bar\mu^tw)=0$; cf. (10.48) and (10.45), respectively. Also, $D^2(r_f) =
D^2(\sigma_P^2) = 0$. Then, the Hessian
